\section{Matrix Operations}
\label{sec:matrix_operations}

\subsection*{MATRIX PROGRAMS 1 --- D1-10.0}
\addcontentsline{toc}{subsection}{MATRIX PROGRAMS 1 --- D1-10.0}
\label{subsec:matrix_programs_1}

\texttt{D1-10.0} is a collection of five routines, all of which utilize the Floating-Point Interpretive System 1 (\texttt{H1-11.0}). This suite consists of the following programs:

\medskip
\begin{tabular}{llll}
\toprule
\textbf{Program} & \textbf{SC Number} & \textbf{Title} \\
\midrule
\texttt{D1-11.0} & 2044               & Matrix Inversion 1 \\
\texttt{D1-12.0} & 2045               & Matrix-Vector Multiplication 1 \\
\texttt{D1-13.0} & 2046               & Matrix Multiplication 1 \\
\texttt{D1-14.0} & 2047               & Matrix Addition/Subtraction 1 \\
\texttt{D1-15.0} & 2048               & Matrix Transposition 1 \\
\bottomrule
\end{tabular}

\subsubsection*{TRANSLATOR NOTE}
Floating-Point Interpreter 1 (\texttt{H1-11.0}, Section~\ref{H1-11.0}) provides a software floating-point implementation for the LGP-21, and is used extensively in the matrix routines.

\newpage

\subsection*{MATRIX INVERSION 1 --- D1-11.0}
\addcontentsline{toc}{subsection}{MATRIX INVERSION 1 --- D1-11.0}
\label{subsec:matrix_inversion_1}

\subsubsection*{PURPOSE}
\texttt{D1-11.0} inverts a square matrix in place, using the Floating-Point Interpretive System 1 (\texttt{H1-11.0}) for its arithmetic. The matrix must be square, non-singular, of rank greater than 1, and stored in floating-point format; it can be as large as available memory allows.

Matrix Inversion 1 is called as a subroutine. Before the call you must exit the Interpretive System (with an \texttt{E} instruction) and patch the interpreter's starting address into the subroutine; the calling sequence then supplies the matrix's rank and starting address. On return, the inverse has replaced the original matrix in memory.

\subsubsection*{STORAGE REQUIREMENTS}
\textbf{Program Storage:} 2 tracks (128 words) \\
\textbf{Temporary Storage:} None

\subsubsection*{INPUT}
The input is the $N^2$ elements of an $N \times N$ matrix in floating-point format, stored row-major (so the first element of row 2 immediately follows the last element of row 1). The matrix must be non-singular.

The subroutine needs an extra column of scratch space, so reserve a contiguous block of $N^2 + N$ locations starting at base address $M$ even though only $N^2$ holds the matrix itself.\footnote{The matrix inversion is implemented using a Gauss-Jordan elimination algorithm that requires a one-column working vector to track pivoting operations and column swaps as it performs the inversion in-place.}

Before calling the subroutine, the starting address of the Interpretive System (see Section~\ref{subsec:complex_interpreter}) must be stored at \texttt{L+163}, where \texttt{L} is the base address of the matrix-inversion code. The operator can do this by hand right after loading the tape, or the main program can do it as part of its setup.

\subsubsection*{CALLING SEQUENCE}
The calling sequence must provide the following parameters:
\begin{enumerate}[label=\arabic*., leftmargin=1.5em, itemsep=2pt]
    \item $N$: The rank (dimension) of the matrix, scaled at $q = 1$, with $N > 1$.
    \item $M$: The starting memory address of the matrix, specified in decimal track/sector notation.
\end{enumerate}
Both parameters $N$ and $M$ are packed into a single parameter word. If $N$ is greater than 15, the packed configuration requires hexadecimal entry format. 

\medskip
\begin{table}[h!]
\centering
\begin{tabularx}{\textwidth}{cccX}
\toprule
\textbf{Memory Location} & \textbf{Instruction} & \textbf{Address} & \textbf{Meaning} \\
\midrule
$\alpha - 1$      & \texttt{E}      & \texttt{0000}    & Exit from the Interpretive System. \\
$\alpha$          & \texttt{R}      & \texttt{L+14}    & Modify subroutine instruction with key word address. \\
$\alpha + 1$      & \texttt{U}      & \texttt{L}       & Unconditional branch to subroutine entry. \\
$\alpha + 2$      & \dots           & \dots            & Packed parameter word: $N$ at $q=15$ and $M$ at $q=29$.\footnote{Yes, you have to work this out yourself. It may be easier in hexadecimal.} \\
$\alpha + 3$      & ---              & ---              & Return to main program. \\
\bottomrule
\end{tabularx}
\end{table}

\medskip
\noindent
The \texttt{E} at $\alpha - 1$ is only needed if the preceding main-program instruction was an interpretive pseudo-instruction; otherwise it can be omitted. The rest of the sequence is the standard \texttt{R}/\texttt{U} call described in Section~\ref{subsec:standard-call2}, with the packed parameter word at $\alpha + 2$ holding $N$ at $q=15$ and $M$ at $q=29$.

\subsubsection*{OUTPUT}
The inverse $B$ replaces the original matrix in memory: $N^2$ row-major elements starting at $M$, same layout as the input.

\subsubsection*{OPERATING DETAILS}

\begin{enumerate}[leftmargin=1.5em, itemsep=6pt]
    \item \textbf{Program Input:} \\
    The subroutine tape is provided in two distribution formats:
    \begin{enumerate}[label=\alph*), itemsep=2pt, topsep=2pt]
        \item arbitrarily relocatable, hexadecimal format, compatible with \texttt{J1-10.0}; and
        \item arbitrarily relocatable, decimal coding sheet format.
    \end{enumerate}
    
    \item \textbf{Required Input/Output Units:} \\
    None.
    
    \item \textbf{Overflow:} \\
    Overflow does not disrupt the entry phase and is not reset upon subroutine exit.
    
    \item \textbf{Error Stops:} \\
    \begin{tabularx}{\textwidth}{lX}
        \toprule
        \textbf{Memory Location} & \textbf{Meaning and Remedy} \\
        \midrule
        \texttt{L+759}           & The matrix is singular. Mathematical inversion is impossible. Execution stops; do not attempt to bypass or resume. \\
        \bottomrule
    \end{tabularx}
\end{enumerate}

\subsubsection*{REMARK}
The subroutine exits the Interpretive System on its own before returning, so control comes back to the main program in native (fixed-point) mode.

\subsubsection*{TRANSLATOR'S NOTES}
Gauss-Jordan elimination is probably the most sophisticated algorithm usable for this purpose on the LGP-21, but it has some pitfalls. Because it performs a larger total number of arithmetic operations than, say, Gaussian elimination with back-substitution, it is even more vulnerable to subtractive cancellation and catastrophic loss of significance.

Subtractive cancellation occurs when you subtract two nearly identical numbers. The leading bits, which match, cancel each other out and leave zeros.

For instance, consider two numbers that match out to 25 bits:
$$\begin{array}{rcc}
1.0110110110110110110110110 & 101 & \times 2^e \\

1.0110110110110110110110110 & 011 & \times 2^e \\
\hline
0.0000000000000000000000000 & 010 & \times 2^e \\
\end{array}$$

After the subtraction, the only remaining non-zero significant bits are the last two. The floating-point interpreter will normalize this result by shifting it left, which shifts in zeroes from the right side of the word. Because the original numbers only had 30 bits of precision, those newly-added lower bits are \emph{garbage}! You now have a value where the error is as large as the signal itself.

In a classic Gauss-Jordan implementation, you aren't just clearing elements below the pivot; you are simultaneously clearing elements \emph{above} the pivot to force the left side of the matrix to be a pure identity matrix. This requires roughly $\frac{n^3}{2}$ multiplications and subtractions compared to the $\frac{n^3}{3}$ required by Gaussian elimination. Every extra row operation is another opportunity for subtractive cancellation to add errors to your data.

If a pivot element becomes very small (usually due to an earlier subtractive cancellation), dividing the rest of the row by that tiny pivot scales up any existing truncation errors dramatically. When that row is subsequently subtracted from all the other rows, that amplified noise propagates through the entire matrix.

Matrix Inversion 1 (Section~\ref{subsec:matrix_inversion_1}) uses a couple of techniques to help prevent this:
\begin{itemize}
	
\item \textbf{Partial Pivoting}: The routine scans down the current column to find the element with the largest absolute magnitude and swaps that row to the pivot position. This ensures that you are always dividing by the largest possible number, keeping the multipliers less than or equal to 1.0, which drastically decreases error amplification.

\item \textbf{The L+759 Hard Stop}: If the largest available pivot is zero (or dangerously close to it under the interpreter's precision bounds), the routine realizes that it has hit a singular or ill-conditioned matrix where subtractive cancellation has wiped out all meaningful data. Rather than returning a useless result, it hits the brakes and halts.
\end{itemize}
\newpage

\subsection*{MATRIX-VECTOR MULTIPLICATION 1 --- D1-12.0}
\addcontentsline{toc}{subsection}{MATRIX-VECTOR MULTIPLICATION 1 --- D1-12.0}
\label{subsec:matrix_vector_multiplication_1}

\subsubsection*{PURPOSE}
\texttt{D1-12.0} multiplies a matrix by a column vector using the Floating-Point Interpretive System 1 (Section~\ref{H1-11.0}), storing the product vector at a specified address. The matrix has $i$ rows and $j$ columns (each less than 64; the matrix doesn't have to be square); the column vector has $j$ elements. All data is in floating-point format, and the total of matrix, vector, and product elements must fit in available memory.

Matrix-Vector Multiplication 1 is called as a subroutine. The Interpretive System must already be loaded; exit it (with an \texttt{E} instruction) before the call. The calling sequence supplies the interpreter's starting address, the matrix's dimensions and starting address, and the starting addresses of the input vector and the output product vector. On return, the product vector occupies the specified locations.

\subsubsection*{STORAGE REQUIREMENTS}
\textbf{Program Storage:} 1 track, 32 sectors (96 words) \\
\textbf{Temporary Storage:} None

\subsubsection*{INPUT}
The inputs are:
\begin{enumerate}[label=\arabic*., leftmargin=1.5em, itemsep=2pt]
    \item an $i \times j$ matrix in row-major order ($i$ and $j$ < 64, not necessarily equal),
    \item a column vector of $j$ elements in consecutive locations.
\end{enumerate}
\noindent
Both must be in Floating-Point Interpretive System 1 format. The matrix, input vector, and output vector together must fit in available memory.

\subsubsection*{CALLING SEQUENCE}
The calling sequence supplies, in order: the Interpretive System's starting address $F$, the matrix's starting address $M$, the matrix's dimensions (row count $i$ as the track part, column count $j$ as the sector part of a single address), the column vector's starting address $V$, and the destination address $P$ for the product vector. All addresses are in decimal track/sector notation. A sample calling sequence (with \texttt{L} the subroutine's base address) is shown below.

\medskip
\begin{table}[h!]
\centering
\begin{tabularx}{\textwidth}{cccX}
\toprule
\textbf{Memory Location} & \textbf{Instruction} & \textbf{Address} & \textbf{Meaning} \\
\midrule
$\alpha - 1$      & \texttt{XE}     & \texttt{0000}    & Exit from the Interpretive System. \\
$\alpha$          & \texttt{R}      & \texttt{$L_0$}   & Store the Interpretive System entry address in subroutine. \\
$\alpha + 1$      & \texttt{U}      & \texttt{$L_0$}   & Unconditional branch to subroutine entry. \\
$\alpha + 2$      & \texttt{Z}      & \texttt{$F_0$}   & Starting address of the Interpretive System. \\
$\alpha + 3$      & \texttt{Z}      & \texttt{M}       & Starting address of the matrix. \\
$\alpha + 4$      & \texttt{Z}      & \texttt{iijj}    & Matrix dimensions: rows (\texttt{ii}) and columns (\texttt{jj}). \\
$\alpha + 5$      & \texttt{Z}      & \texttt{V}       & Starting address of the column vector. \\
$\alpha + 6$      & \texttt{Z}      & \texttt{P}       & Starting address of the destination product vector. \\
$\alpha + 7$      & ---              & ---              & Return to main program. \\
\bottomrule
\end{tabularx}
\end{table}

\medskip
\noindent
The \texttt{E} at $\alpha-1$ is only needed if the preceding main-program instruction was an interpretive pseudo-instruction. The rest is the standard inlined-parameter call described in Section~\ref{subsec:standard-call2}; locations $\alpha+2$ through $\alpha+6$ are the parameter list, and $\alpha+7$ is where execution resumes after the call.

\subsubsection*{OUTPUT}
The product vector is written to consecutive locations starting at $P$.

\subsubsection*{OPERATING DETAILS}

\begin{enumerate}[leftmargin=1.5em, itemsep=6pt]
    \item \textbf{Program Input:} \\
    The subroutine tape is distributed in two formats:
    \begin{enumerate}[label=\alph*), itemsep=2pt, topsep=2pt]
        \item arbitrarily relocatable, hexadecimal format, compatible with \texttt{J1-10.0}; and
        \item arbitrarily relocatable, decimal coding sheet format.
    \end{enumerate}
    
    \item \textbf{Required Input/Output Units:} \\
    None.
    
    \item \textbf{Overflow:} \\
    Overflow is ignored during the execution initialization phase and is not reset upon subroutine exit.
\end{enumerate}

\subsubsection*{NOTES}
The subroutine exits the Interpretive System on its own before returning, so control comes back to the main program in native (fixed-point) mode.

\subsubsection*{TRANSLATOR'S NOTE}
This is the first function to use the extended calling sequence that looks like this, documented at Section~\ref{subsec:standard-call2}.

Speculating slightly, because I have not seen/dumped the code in question, this is my best guess on how it works.
\begin{itemize}
	\item The \texttt{R} instruction in the main program stores the return address ($\alpha+2$) into $L_0$. Note that the current address + 2 is the start of the parameter list to be passed to the subroutine.
	\item The \texttt{U} instruction at $\alpha+1$ jumps to $L_0$.
	\item A \texttt{B} instruction at $L_0$ loads $\alpha+2$, the parameter list's address into the accumulator.
	\item The subroutine adds 5 to the accumulator to get the actual return address (the word following the parameter list).
	\item The subroutine then stores the return address into a \texttt{U} instruction that it will use for the return to the main program.
\end{itemize}
\noindent
This lets the \texttt{B} instruction at $L_0$ do double duty: it both stashes the parameter-list address (which the subroutine treats as a stand-in for the return address) and loads that address into the accumulator so the subroutine can compute the real return point.
\newpage

\subsection*{MATRIX MULTIPLICATION 1 --- D1-13.0}
\addcontentsline{toc}{subsection}{MATRIX MULTIPLICATION 1 --- D1-13.0}
\label{subsec:matrix_multiplication_1}

\subsubsection*{PURPOSE}
\texttt{D1-13.0} computes the matrix product $C = AB$ using the Floating-Point Interpretive System 1 (\texttt{H1-11.0}), where $A$ is $n \times m$, $B$ is $m \times p$, and $C$ is $n \times p$. All inputs are in floating-point format. The matrix dimensions must each be less than 64.\footnote{\textbf{Translator's note:} The limit of 64 is an architectural artifact of the LGP-21's memory, which is physically structured as 64 tracks of 64 sectors each. Capping dimensions at 64 allows index calculation loops to stay bound within simple track-selection logic without risking multi-track pointer overflow.} The total of all elements of $A$, $B$, and $C$ must fit in available memory.

Matrix Multiplication 1 is called as a subroutine. Exit the Interpretive System (with an \texttt{E} instruction) before the call. The calling sequence supplies the interpreter's starting address, the dimensions and starting addresses of $A$ and $B$, and the starting address of $C$. On return, the product matrix occupies the destination locations.

\subsubsection*{STORAGE REQUIREMENTS}
\textbf{Program Storage:} 2 tracks, 50 sectors (178 words) \\
\textbf{Temporary Storage:} None

\subsubsection*{INPUT}
The inputs are matrix $A$ ($n \times m$) and matrix $B$ ($m \times p$), both in row-major order, both in Floating-Point Interpretive System 1 format (Section~\ref{H1-11.0}). The two matrices plus the destination $C$ require $(n \times m) + (m \times p) + (n \times p)$ contiguous locations of memory in total.

\subsubsection*{CALLING SEQUENCE}
The calling sequence supplies, in order: the Interpretive System's starting address $F$, the starting address of matrix $A$, the dimensions of $A$ (rows $n$ as the track part, columns $m$ as the sector part), the starting address of matrix $B$, the dimensions of $B$ (rows $m$ as the track part, columns $p$ as the sector part), and the starting address of the destination matrix $C$. All addresses are in decimal track/sector notation. A template calling sequence (with \texttt{$L_0$} the subroutine's base address) is shown below.

\medskip
\begin{table}[h!]
\centering
\begin{tabularx}{\textwidth}{cccX}
\toprule
\textbf{Memory Location} & \textbf{Instruction} & \textbf{Address} & \textbf{Meaning} \\
\midrule
$\alpha - 1$      & \texttt{E}      & \texttt{0000}    & Exit from the Interpretive System. \\
$\alpha$          & \texttt{R}      & \texttt{$L_0$}   & Store parameter keyword tracking address in subroutine.\footnote{\textbf{Translator's note:} As discussed at Section~\ref{subsec:standard-call2}, the \texttt{R} calculates the location $\alpha+2$, effectively the parameter list address, and embeds it into the instruction at \texttt{$L_0$}. The jump to the subroutine executes a \texttt{B} which loads the address of the parameters, adds an offset of $+8$ to compute the final return address, and saves it locally to properly do the return.} \\
$\alpha + 1$      & \texttt{U}      & \texttt{$L_0$}   & Unconditional branch to subroutine entry. \\
$\alpha + 2$      & \texttt{Z}      & \texttt{$F_0$}   & Starting address of the Interpretive System. \\
$\alpha + 3$      & \texttt{Z}      & \texttt{A}       & Starting address of matrix $A$. \\
$\alpha + 4$      & \texttt{Z}      & \texttt{nnmm}    & Dimensions of matrix $A$: rows (\texttt{nn}) and columns (\texttt{mm}). \\
$\alpha + 5$      & \texttt{Z}      & \texttt{B}       & Starting address of matrix $B$. \\
$\alpha + 6$      & \texttt{Z}      & \texttt{mmpp}    & Dimensions of matrix $B$: rows (\texttt{mm}) and columns (\texttt{pp}). \\
$\alpha + 7$      & \texttt{Z}      & \texttt{C}       & Starting address of destination product matrix $C$. \\
$\alpha + 8$      & ---              & ---              & Return to main program (execution resumes here). \\
\bottomrule
\end{tabularx}
\end{table}

\medskip
\noindent
The \texttt{E} at $\alpha-1$ is only needed if the preceding main-program instruction was an interpretive pseudo-instruction.

\subsubsection*{OUTPUT}
$C$ is written in row-major order to consecutive locations starting at $C$, $n \times p$ elements total.

\subsubsection*{OPERATING DETAILS}

\begin{enumerate}[leftmargin=1.5em, itemsep=6pt]
    \item \textbf{Program Input:} \\
    The subroutine tape is distributed in two formats:
    \begin{enumerate}[label=\alph*), itemsep=2pt, topsep=2pt]
        \item arbitrarily relocatable, hexadecimal format, compatible with \texttt{J1-10.0}; and
        \item arbitrarily relocatable, decimal coding sheet format.
    \end{enumerate}
    
    \item \textbf{Required Input/Output Units:} \\
    None.
    
    \item \textbf{Overflow:} \\
    Overflow is ignored during the entry configuration phase and is not reset upon subroutine exit.
    
    \item \textbf{Error Stops:} \\
    \begin{tabularx}{\textwidth}{lX}
        \toprule
        \textbf{Memory Location} & \textbf{Meaning and Remedy} \\
        \midrule
        \texttt{L+241}           & Inner matrix dimensions mismatch: the number of columns in matrix $A$ ($m$) does not equal the number of rows in matrix $B$ ($m$). Mathematical multiplication is undefined. Do not attempt to resume execution; the input data dimensions must be corrected. \\
        \bottomrule
    \end{tabularx}
\end{enumerate}

\subsubsection*{NOTES}
The subroutine exits the Interpretive System on its own before returning, so control comes back to the main program in native (fixed-point) mode.
\newpage

\subsection*{MATRIX ADDITION AND SUBTRACTION 1 --- D1-14.0}
\addcontentsline{toc}{subsection}{MATRIX ADDITION AND SUBTRACTION 1 --- D1-14.0}
\label{subsec:matrix_addition_subtraction_1}

\subsubsection*{PURPOSE}
\texttt{D1-14.0} adds or subtracts two $i \times j$ matrices $A$ and $B$ using the Floating-Point Interpretive System 1 (\texttt{H1-11.0}), storing the result in a third memory area. The matrices don't need to be square; $i$ and $j$ must each be less than 64. The two inputs and the destination must all fit in available memory.

Matrix Addition and Subtraction 1 is called as a subroutine. Exit the Interpretive System (with an \texttt{E} instruction) before the call. The calling sequence supplies the interpreter's starting address, the matrix dimensions and starting addresses, the desired operation (add or subtract), and the destination address. On return, the result occupies the destination.

\subsubsection*{STORAGE REQUIREMENTS}
\textbf{Program Space:} 1 track (64 words) \\
\textbf{Temporary Storage:} None

\subsubsection*{INPUT}
The inputs are matrices $A$ and $B$, both $i \times j$ and in identical layouts --- both row-major or both column-major.\footnote{This is significantly more flexible than the Matrix Multiplication routine (\texttt{D1-13.0}). Because addition and subtraction are purely element-by-element operations, the subroutine doesn't actually care about row vs. column orientation, as long as both inputs share the exact same structure. The execution loop simply scans linearly across both matrices.} Both must be in Floating-Point Interpretive System 1 format, with $i$ and $j$ each less than 64. The matrices and the destination together require $3(ij)$ locations of memory in total.

\subsubsection*{CALLING SEQUENCE}
The calling sequence supplies, in order: the Interpretive System's starting address $F_0$, the starting address of matrix $A$, the matrix dimensions (rows $i$ as the track part, columns $j$ as the sector part), the operation word --- an \texttt{A} or \texttt{S} opcode with $B$'s starting address as its operand\footnote{Another memory-saving optimization: instead of wasting an entire word on an operation flag or setting up a parsing branch, the subroutine reads this parameter word at $\alpha+5$ and injects it directly into its internal math loop as an executable instruction; the address of the second matrix $B$ is \emph{also} packed into this instruction; the instruction (\texttt{A} for Add, \texttt{S} for Subtract) is executed dynamically against the current matrix address via self-modifying code.} --- and the destination address $C$. All addresses are in decimal track/sector notation. A template calling sequence (with \texttt{$L_0$} the subroutine's base address) is shown below.

\medskip
\begin{table}[h!]
\centering
\begin{tabularx}{\textwidth}{cccX}
\toprule
\textbf{Memory Location} & \textbf{Instruction} & \textbf{Address} & \textbf{Meaning} \\
\midrule
$\alpha - 1$      & \texttt{E}      & \texttt{0000}    & Exit from the Interpretive System. \\
$\alpha$          & \texttt{R}      & \texttt{$L_0$}   & Store parameter tracking address in the subroutine. \\
$\alpha + 1$      & \texttt{U}      & \texttt{$L_0$}   & Unconditional branch to subroutine entry. \\
$\alpha + 2$      & \texttt{Z}      & \texttt{$F_0$}   & Starting address of the Interpretive System. \\
$\alpha + 3$      & \texttt{Z}      & \texttt{A}       & Starting address of matrix $A$. \\
$\alpha + 4$      & \texttt{Z}      & \texttt{iijj}    & Dimensions of the matrices: rows (\texttt{ii}) and columns (\texttt{jj}). \\
$\alpha + 5$      & \texttt{A}/\texttt{S} & \texttt{B}  & Operation (\texttt{A} or \texttt{S}) and starting address of matrix $B$. \\
$\alpha + 6$      & \texttt{Z}      & \texttt{C}       & Starting address of destination area for the resulting matrix. \\
$\alpha + 7$      & ---              & ---              & Return to main program (execution resumes here). \\
\bottomrule
\end{tabularx}
\end{table}

\medskip
\noindent
The \texttt{E} at $\alpha-1$ is only needed if the preceding main-program instruction was an interpretive pseudo-instruction.

\subsubsection*{OUTPUT}
The result is written in floating-point format to consecutive locations starting at $C$, $ij$ elements total.

\subsubsection*{OPERATING DETAILS}

\begin{enumerate}[leftmargin=1.5em, itemsep=6pt]
    \item \textbf{Program Input:} \\
    The subroutine tape is distributed in two formats:
    \begin{enumerate}[label=\alph*), itemsep=2pt, topsep=2pt]
        \item arbitrarily relocatable, hexadecimal format, compatible with \texttt{J1-10.0}; and
        \item arbitrarily relocatable, decimal coding sheet format.
    \end{enumerate}
    
    \item \textbf{Required Input/Output Units:} \\
    None.
    
    \item \textbf{Overflow:} \\
    Overflow is ignored during the entry initialization phase and is not reset upon subroutine exit.
\end{enumerate}

\subsubsection*{REMARKS}
The subroutine exits the Interpretive System on its own before returning, so control comes back to the main program in native (fixed-point) mode.

The destination can be matrix $A$, matrix $B$, or a separate region. If you use a separate region, its memory range must not overlap with $A$ or $B$.
\newpage

\subsection*{MATRIX TRANSPOSITION 1 --- D1-15.0}
\addcontentsline{toc}{subsection}{MATRIX TRANSPOSITION 1 --- D1-15.0}
\label{subsec:matrix_transposition_1}

\subsubsection*{PURPOSE}
\texttt{D1-15.0} transposes a square matrix in memory, writing the result to a destination address (which may be the same as the source). The matrix's rank must be less than 32. Both fixed-point and floating-point formats are supported. The matrix and (if separate) the destination together must fit in available memory.

Matrix Transposition 1 is called as a subroutine. Exit the Interpretive System (with an \texttt{E} instruction) before the call. The calling sequence supplies the matrix's rank and starting address and the destination address. On return, the transposed matrix occupies the destination.

\subsubsection*{STORAGE REQUIREMENTS}
\textbf{Program Space:} 1 track, 58 sectors (122 words) \\
\textbf{Temporary Storage:} None

\subsubsection*{INPUT}
The input is the elements of a square matrix in consecutive locations, either row-major or column-major. The data can be in raw fixed-point fractional format or in Floating-Point Interpretive System 1 format (\texttt{H1-11.0}). The rank $N$ must satisfy $2 < N \le 31$. The matrix takes $N^2$ locations; if the destination is separate, total memory required is $2N^2$.

\subsubsection*{CALLING SEQUENCE}
The calling sequence supplies, in order: the matrix's rank $N$ (at $q=15$) and starting address $M$ --- packed into a single parameter word at $\alpha+2$, requiring hexadecimal entry if $N > 15$ --- and the destination address $T$. All addresses are in decimal track/sector notation. $T$ may equal $M$ (in-place transposition); otherwise, the destination region must not overlap the source. A template calling sequence (with \texttt{L} the subroutine's base address) is shown below.

\medskip
\begin{table}[h!]
\centering
\begin{tabularx}{\textwidth}{cccX}
\toprule
\textbf{Memory Location} & \textbf{Instruction} & \textbf{Address} & \textbf{Meaning} \\
\midrule
$\alpha - 1$      & \texttt{E}      & \texttt{0000}    & Exit from the Interpretive System. \\
$\alpha$          & \texttt{R}      & \texttt{L}       & Store the parameter keyword tracking address in the subroutine. \\
$\alpha + 1$      & \texttt{U}      & \texttt{L}       & Unconditional branch to subroutine entry. \\
$\alpha + 2$      & \dots           & \dots            & Packed parameter word: $N$ at $q=15$ and $M$ at $q=29$. \\
$\alpha + 3$      & \texttt{Z}      & \texttt{T}       & Starting address for the transposed matrix. \\
$\alpha + 4$      & ---              & ---              & Return to main program (execution resumes here). \\
\bottomrule
\end{tabularx}
\end{table}

\medskip
\noindent
The \texttt{E} at $\alpha-1$ is only needed if the preceding main-program instruction was an interpretive pseudo-instruction.

\subsubsection*{OUTPUT}
The transposed matrix is written to consecutive locations starting at $T$, in the same row/column orientation as the input (so a row-major input produces a row-major output of the transpose, and likewise for column-major).

\subsubsection*{OPERATING DETAILS}

\begin{enumerate}[leftmargin=1.5em, itemsep=6pt]
    \item \textbf{Program Input:} \\
    The subroutine tape is distributed in two formats:
    \begin{enumerate}[label=\alph*), itemsep=2pt, topsep=2pt]
        \item arbitrarily relocatable, hexadecimal format, compatible with \texttt{J1-10.0}; and
        \item arbitrarily relocatable, decimal coding sheet format.
    \end{enumerate}
    
    \item \textbf{Required Input/Output Units:} \\
    None.
    
    \item \textbf{Overflow:} \\
    Overflow is ignored during the execution initialization phase and is not reset upon subroutine exit.
\end{enumerate}

\subsubsection*{REMARK}
Unlike the other routines in the matrix suite, this subroutine \emph{does not} invoke or use the Floating-Point Interpretive System internally.\footnote{Because matrix transposition is a purely structural index manipulation (swapping element positions $(i,j)$ with $(j,i)$), the subroutine does not perform any actual arithmetic on the matrix values themselves, so it's completely unnecessary to incur the massive overhead of the floating-point interpreter \texttt{H1-11.0}. We're just moving data, so it doesn't matter if it's native fixed-point integers or multi-word floating-point interpreter variables.}
\newpage
%%%%%%%%%%%%%%%%%%%%%%%%%%%%%%%%%%%%%%%%%%%%%%%%%%%%%%%%%%%%%%%%%%%%%%%%%%%%%%%%%
\subsection*{COMPLEX OPERATIONS INTERPRETIVE SYSTEM --- H1-10.0}
\addcontentsline{toc}{subsection}{COMPLEX OPERATIONS INTERPRETIVE SYSTEM --- H1-10.0}
\label{subsec:complex_interpreter}

\subsubsection*{PURPOSE}
\texttt{H1-10.0} enables the user to process algebraic expressions involving complex numbers in the form $(a+bi)$ as seamlessly as if they were standard real numbers. Additionally, the system provides integrated left and right logical bit-shifting capabilities. It also allows for dynamic address modification and loop-termination testing against a target end-address without ever needing to exit the interpretive execution loop.

\subsubsection*{STORAGE REQUIREMENTS}
\textbf{Program Space:} 3 tracks (192 words) \\
\textbf{Temporary Storage:} None

\subsubsection*{INPUT}
For every complex data element, both the real and imaginary parts must be stored at the identical binary scaling factor ($q$-factor) in consecutive locations. Specifically, the real part must reside at location \texttt{ttss} and the imaginary part at location \texttt{ttss+1}.

\noindent
Because each complex number occupies two discrete memory locations, the subroutine implements simulated registers to facilitate instruction execution. These architecture registers consist of a simulated pseudo-accumulator located at location \texttt{L+33} and a simulated address register located at location \texttt{L+219}, where \texttt{L} denotes the starting base address of the subroutine. In addition to these active components, the program records the transient state of the physical hardware accumulator at location \texttt{L+59}.

\noindent
The interpretive engine processes 15 distinct operations. For all arithmetic commands, the address field \texttt{ttss} references the location containing the real component, though the internal loop automatically processes both parts of the complex value. For all other non-arithmetic commands, \texttt{ttss} functions as a standard single-word memory address. Only the instructions explicitly detailed below are structurally valid within an interpretive instruction sequence.

\paragraph{1. Arithmetic Instructions}
\begin{description}[leftmargin=1.5em, itemsep=4pt]
    \item[\texttt{Bttss} (Bring):] The elements at locations \texttt{ttss} and \texttt{ttss+1} replace the contents of the pseudo-accumulator. The source memory remains unchanged.
    \item[\texttt{Attss} (Add):] The elements at locations \texttt{ttss} and \texttt{ttss+1} are added to the pseudo-accumulator. The result replaces the contents of the pseudo-accumulator.
    \item[\texttt{Sttss} (Subtract):] The elements at locations \texttt{ttss} and \texttt{ttss+1} are subtracted from the pseudo-accumulator. The result replaces the contents of the pseudo-accumulator.
    \item[\texttt{Mttss} (Multiply):] The pseudo-accumulator is multiplied by the complex value stored at \texttt{ttss} and \texttt{ttss+1}. The product replaces the contents of the pseudo-accumulator.\footnote{\textbf{Translator's note:} A complex multiplication $(a+bi)(c+di)$ requires computing $(ac-bd) + (ad+bc)i$. Performing this in a single interpretive step saves a massive amount of main-program setup and temporary storage management.}
    \item[\texttt{Dttss} (Divide):] The pseudo-accumulator is divided by the complex value stored at \texttt{ttss} and \texttt{ttss+1}. The quotient replaces the contents of the pseudo-accumulator.
    \item[\texttt{Httss} (Hold):] The contents of the pseudo-accumulator are copied into locations \texttt{ttss} and \texttt{ttss+1}. The pseudo-accumulator remains unchanged.
    \item[\texttt{Cttss} (Clear):] The contents of the pseudo-accumulator are copied into locations \texttt{ttss} and \texttt{ttss+1}, and the pseudo-accumulator is subsequently reset to zero.
\end{description}

\paragraph{2. Auxiliary Instructions}
\begin{description}[leftmargin=1.5em, itemsep=4pt]
    \item[\texttt{XRO0nn} (Shift Right):] Shifts both components of the pseudo-accumulator right by \texttt{nn} binary places ($0 \le \texttt{nn} \le 10$).
    \item[\texttt{XPO0nn} (Shift Left):] Shifts both components of the pseudo-accumulator left by \texttt{nn} binary places ($0 \le \texttt{nn} \le 10$).
    \item[\texttt{Uttss} (Unconditional Jump):] Transfers execution control. The next pseudo-instruction to be interpreted is fetched from location \texttt{ttss}. This jump cannot be used to exit the interpreter; \texttt{ttss} must point to a valid complex pseudo-instruction.
    \item[\texttt{XE0000} (Exit):] Explicit exit from the Interpretive System. The instruction immediately following \texttt{XE0000} is executed as a native hardware instruction.
\end{description}

\paragraph{3. Address Modification Instructions}
\begin{description}[leftmargin=1.5em, itemsep=4pt]
    \item[\texttt{Ettss} (Set Address Register):] The address portion of the word stored at location \texttt{ttss} is loaded into the simulated address register. Memory remains unchanged.
    \item[\texttt{Ittss} (Increment Address):] Increments the current value of the simulated address register by the value \texttt{ttss}.
    \item[\texttt{Yttss} (Inject Address):] Embeds the current value of the simulated address register directly into the address portion of the instruction word at location \texttt{ttss}. The address register itself and all other memory remains unchanged.
    \item[\texttt{Zttss} (Test and Skip):] Subtracts the current value of the simulated address register from the literal value \texttt{ttss}. If the result is exactly zero, the next pseudo-instruction is skipped (execution jumps to the second following word). If the result is non-zero, execution continues normally with the immediately following instruction.
\end{description}

\subsubsection*{CALLING SEQUENCE}
A template for entering the interpretive loop is illustrated below, where \texttt{L} denotes the starting address of the interpreter.

\medskip
\begin{table}[h!]
\centering
\begin{tabularx}{\textwidth}{cccX}
\toprule
\textbf{Memory Location} & \textbf{Instruction} & \textbf{Address} & \textbf{Meaning} \\
\midrule
$\alpha$          & \texttt{R}      & \texttt{L}       & Store parameter tracking address inside the interpreter. \\
$\alpha + 1$      & \texttt{U}      & \texttt{L}       & Unconditional branch to interpreter entry. \\
$\alpha + 2$      & \dots           & \dots            & First complex pseudo-instruction to be interpreted. \\
$\vdots$          & $\vdots$        & $\vdots$         & Sequence of complex operations. \\
$\alpha + n$      & \texttt{XE}     & \texttt{0000}    & Exit from the Interpretive System. \\
$\alpha + n + 1$  & \dots           & \dots            & First native hardware instruction (execution resumes here). \\
\bottomrule
\end{tabularx}
\end{table}

\medskip
\noindent
The \texttt{R} instruction at location $\alpha$ automatically captures the address of the initial pseudo-instruction ($\alpha + 2$) and saves it internally within the interpreter's program counter tracker.

\subsubsection*{OPERATING DETAILS}

\paragraph{1. Execution Time}
The following table outlines the maximum execution durations required for each individual complex operation, measured in drum rotations and milliseconds.\footnote{Notice how left-shifting (the \texttt{P} instruction) is significantly slower than right-shifting (\texttt{R}). On a magnetic disk memory system like the LGP-21, bit-shifting requires track-recirculation delays and masking. Left shifts require repeated additions or multiplication loops to slide data up toward the sign bit, adding extra revolutions, whereas right shifts could take advantage of standard hardware truncation.}

\medskip
\noindent
\small
\begin{tabular}{crr|crr}
\toprule
\textbf{Instruction} & \textbf{Rotations} & \textbf{Time (ms)} & \textbf{Instruction} & \textbf{Rotations} & \textbf{Time (ms)} \\
\midrule
\texttt{B}           & 11 & 561  & \texttt{P}            & 17 & 1167 \\
\texttt{A}           & 14 & 714  & \texttt{U}            & 8  & 408  \\
\texttt{S}           & 14 & 714  & \texttt{XE} (Exit)    & 6  & 306  \\
\texttt{M}           & 22 & 1122 & \texttt{E} (Set Addr) & 9  & 459  \\
\texttt{D}           & 41 & 2091 & \texttt{I}            & 6  & 306  \\
\texttt{H}           & 14 & 714  & \texttt{Y}            & 8  & 408  \\
\texttt{C}           & 14 & 714  & \texttt{Z} ($=0$)     & 11 & 561  \\
\texttt{R} (Shift R) & 14 & 714  & \texttt{Z} ($\ne0$)   & 13 & 663  \\
\bottomrule
\end{tabular}
\normalsize

\paragraph{2. Program Input:}
The subroutine tape is distributed in two formats:
\begin{enumerate}[label=\alph*), itemsep=2pt, topsep=2pt]
    \item arbitrarily relocatable, hexadecimal format, compatible with \texttt{J1-10.0}; and
    \item arbitrarily relocatable, decimal coding sheet format.
\end{enumerate}

\paragraph{3. Required Input/Output Units:} None.

\paragraph{4. Error Stops:}
\begin{center}
\begin{tabular}{l p{6.5cm} l}
    \toprule
    \textbf{Location} & \textbf{Meaning} & \textbf{Remedy} \\
    \midrule
    \texttt{L+122}    & Negative Instruction encountered (Illegal format). & Fatal error. \\
    \texttt{L+135}    & Invalid \texttt{N-} Type instruction variation.    & Fatal error. \\
    \texttt{L+154}    & Invalid \texttt{T-} Type instruction variation.    & Fatal error. \\
    \bottomrule
\end{tabular}
\end{center}

\paragraph{5. Overflow:}
Overflow status is ignored during interpretive entry execution loops and is not explicitly reset upon subroutine exit.

\subsubsection*{REMARKS}
The first instruction inside the interpretive sequence (location $\alpha+2$) must not be an \texttt{XE0000} command, as it would cause an immediate, unintended system exit before any operation can occur.

\subsubsection*{TRANSLATOR NOTES}
The inclusion of index/address pseudoinstructions (\texttt{E}, \texttt{I}, \texttt{Y}, \texttt{Z}) elevates this interpreter from a simple complex-math math library to a self-contained environment. It can handle full multi-row matrix array traversal directly inside the complex domain without ever exiting the interpreter loop to do data manipulation and control flow.
\newpage
